\documentclass[11pt,a4paper]{article}
\usepackage[T1]{fontenc}
\usepackage[utf8]{inputenc}
\usepackage[brazil]{babel}
\usepackage{lmodern}
\usepackage[margin=2.5cm]{geometry}
\usepackage{xcolor}
\usepackage[hidelinks]{hyperref}
\usepackage{microtype}
\setlength{\parskip}{0.6em}
\setlength{\parindent}{0pt}
\setlength{\emergencystretch}{3em}
\newcommand{\campo}[1]{\texttt{#1}}

\title{\textbf{Levantamento das bases de dados}\\[0.3em]\large Projeto Extensionista: análise de dados de enchentes do INCT SiMAI\\Coleta, Preparação e Análise de Dados (Fase 1)}
\author{Autor: João Pedro da Silva Kriger}
\date{21 de novembro de 2026}

\begin{document}
\maketitle

\section{Objetivo}

Este relatório descreve as bases sobre enchentes no Rio Grande do Sul levantadas para o projeto: de onde vêm, o que trazem, que recorte de espaço e tempo cobrem, em que formato estão, como foram coletadas, sob que condições podem ser usadas e como se ligam aos dados do INCT SiMAI. No fim, justifica a escolha das bases da Tela 2 e da pergunta que ela responde.

O levantamento começou em 27/10/2026. Esta versão é a reentrega da Fase 1 e descreve as bases que de fato ficaram no projeto; as que foram avaliadas e descartadas estão na seção~\ref{sec:descartadas}. A origem de cada arquivo também está em \campo{dados/README.md}, no repositório.

\section{Simulação do INCT SiMAI}

\textbf{Origem.} Instituto Nacional de Ciência e Tecnologia SiMAI, que publica o resultado na página \href{http://vhlab.com.br/INCT/floods/}{Flood -- INCT SiMAI}.

\textbf{Conteúdo.} O arquivo \campo{data.json} traz, para cada nível do Guaíba, os 94 bairros de Porto Alegre com cinco campos: total de prédios (OpenStreetMap), prédios atingidos (simulação), porcentagem de prédios atingidos, população do bairro (Censo IBGE 2022) e população afetada estimada. Além dele, foram coletadas as manchas simuladas, ou seja, a área que a simulação indica como alagada em cada nível.

\textbf{Recorte.} Município de Porto Alegre. Níveis de 3,0 a 10,0 m em passos de 0,1 m (71 níveis, 6.674 registros); as manchas foram coletadas de 3,0 a 5,4 m (25 níveis). Não há dimensão de tempo: é um cenário por nível do rio.

\textbf{Formato.} JSON aninhado (nível, bairro, campos). As manchas são polígonos em JSON.

\textbf{Coleta.} O \campo{data.json} foi fornecido pela disciplina em 27/10/2026. As manchas vêm da mesma API usada pelo site do INCT (\campo{POST http://vhlab.lad.pucrs.br:5002/api/map-data}), coletadas por código em 10/11/2026 com \campo{scripts/coleta\_bases.py --temas inct}. A API não tem \campo{robots.txt}, e o script faz pausas entre as requisições.

\textbf{Licença.} Material cedido para fins acadêmicos no projeto; o uso cita o INCT SiMAI como fonte.

\textbf{Ligação com o INCT.} É a própria simulação: o JSON é a única base da Tela 1, e as manchas ligam a Tela 2 ao mesmo nível do rio.

\section{Malha oficial de bairros (Prefeitura de Porto Alegre)}

\textbf{Origem.} Unidade de Monitoramento do Desenvolvimento Urbano, Secretaria Municipal de Urbanismo, no servidor ArcGIS da Prefeitura (camada \campo{A02\_SOLO\_CRIADO/bairros}). O enunciado indica o ObservaPOA; foi usada a camada oficial publicada pela própria Prefeitura no seu servidor de mapas, com os limites definidos em lei.

\textbf{Conteúdo.} Os limites dos 94 bairros definidos pela Lei Municipal 12.112/2016.

\textbf{Recorte.} Município de Porto Alegre; limites vigentes desde 2016.

\textbf{Formato.} GeoJSON, em versão simplificada para exibição.

\textbf{Coleta.} Por código, pela API REST do ArcGIS, em 27/10/2026 (\campo{scripts/coleta\_bases.py --temas arcgis}). O servidor não tem \campo{robots.txt}; o script faz pausas.

\textbf{Licença.} Dado aberto publicado pela Prefeitura, sem licença explícita na camada; o uso cita a fonte.

\textbf{Ligação com o INCT.} Pelo nome do bairro. Dois nomes são escritos de forma diferente nas duas fontes, e a equivalência ficou na coluna \campo{bairro\_malha} do dado limpo. A malha é usada só para desenhar o mapa da Tela 1 e as divisas na Tela 2.

\section{Monitoramento da Reconstrução (Prefeitura de Porto Alegre)}

\textbf{Origem.} Dados abertos do painel \href{https://prefeitura.poa.br/monitoramento-reconstrucao}{Monitoramento da Reconstrução}, publicados no servidor ArcGIS da Prefeitura (serviços \campo{HUB/base\_reconstrucao}, \campo{HUB/hub\_habitacao\_} e \campo{HUB/hub\_base\_financeiro}).

\textbf{Conteúdo.} As 318 obras de reconstrução que a Prefeitura monitora desde a enchente de 2024 (194 de equipamentos públicos e 124 do sistema de proteção contra cheias), com tipo, bairro, coordenada, valor previsto, fonte do recurso e situação; as 570 etapas das obras do sistema de proteção, com contratos e empresas; as 120 áreas de habitação do DEMHAB, com a necessidade estimada de reposição; e o painel financeiro, com a necessidade estimada, os valores solicitados à União e recebidos, e os recursos aplicados por origem.

\textbf{Recorte.} Município de Porto Alegre, obras e recursos posteriores à enchente de maio de 2024. Os registros são de 17/09/2026; a base não informa a data de cada situação de obra.

\textbf{Formato.} JSON da API do ArcGIS (atributos e geometria).

\textbf{Coleta.} Por código, pela API REST do ArcGIS, em 27/10/2026, com conferência em 09/11/2026 de que o número de registros é igual ao da fonte.

\textbf{Licença.} Dado aberto publicado pela Prefeitura, sem licença explícita nas camadas; o uso cita o painel e a data dos registros.

\textbf{Ligação com o INCT.} Pela localização: cada obra com coordenada é cruzada com a mancha do INCT em cada nível do rio, e pelo bairro.

\section{Contratos e empenhos (TCE-RS)}

\textbf{Origem.} \href{https://dados.tce.rs.gov.br}{Portal de dados abertos do Tribunal de Contas do Estado do RS}.

\textbf{Conteúdo.} Licitações, contratos (com aditivos) e empenhos da Prefeitura (código 54900) e do DMAE (54907).

\textbf{Recorte.} Porto Alegre, anos de 2024 e 2025.

\textbf{Formato.} CSV compactado em ZIP.

\textbf{Coleta.} Por código, só por URLs diretas de download e com 10 segundos entre acessos, porque o \campo{robots.txt} do portal proíbe o caminho \campo{/api/} e pede esse intervalo (\campo{scripts/coleta\_bases.py --temas tce}).

\textbf{Licença e cuidado.} Dados abertos do TCE-RS, sem licença explícita nos conjuntos consultados; o uso cita a fonte. Os arquivos contêm CPF e nome de credores pessoa física, por isso não são versionados no repositório e nenhum nome de pessoa física é publicado.

\textbf{Ligação com o INCT.} Nenhuma direta. A base serviu só para conferir os valores das obras do DMAE e não aparece no dashboard.

\section{Prédios atingidos em 2024 (Copernicus EMSN194)}

\textbf{Origem.} \href{https://riskandrecovery.emergency.copernicus.eu/EMSN194/}{Copernicus Emergency Management Service}, ativação EMSN194, produto P08.

\textbf{Conteúdo.} 119.626 prédios e instalações marcados como atingidos por intérpretes em imagens de 0,5 m, dos quais 42,6 mil em Porto Alegre.

\textbf{Recorte.} Áreas de interesse 02, 03 e 04 da ativação, que cobrem só parte da cidade; imagens de maio de 2024.

\textbf{Formato.} GeoJSON de pontos.

\textbf{Coleta.} Manual, porque o \campo{robots.txt} do site proíbe o caminho \campo{/download/}. O script de coleta indica o link e a pasta. Baixado em 04/11/2026.

\textbf{Licença.} Uso livre com a citação ``© European Union, Copernicus Emergency Management Service''.

\textbf{Ligação com o INCT.} Pela enchente de 2024: os pontos aparecem como camada opcional no mapa da Tela 2, para comparar a área simulada com os locais atingidos de verdade.

\section{Bases avaliadas e descartadas}
\label{sec:descartadas}

Foram avaliadas e retiradas: a mancha observada do IPH/UFRGS e os domicílios do Censo, para validar a mancha do INCT, porque não existe uma fonte da verdade para a extensão da enchente; o nível do Guaíba da ANA e as áreas das casas de bombas, para estudar falhas do sistema de proteção, retiradas para manter o foco; a fila do SUS (GERCON), que só tinha consultas com especialista; e as empresas sancionadas, CNPJ e licitações do Dados Abertos POA, que serviam a uma busca de pontos de atenção em contratos que deixou de ser o foco. Os dados dessas análises foram apagados.

\section{Escolha das bases e da pergunta da Tela 2}

O enunciado pede que a Tela 2 mostre o que outras bases acrescentam à resposta da Tela 1. A Tela 1 diz quem a água alcança; a pergunta escolhida para a Tela 2 é: \textbf{se o rio subir, quais obras da reconstrução ficam na área alcançada pela água, e quanto está previsto nelas?}

A base central é o Monitoramento da Reconstrução da Prefeitura, por três motivos. Ela é oficial, aberta e obtida por código. Ela tem coordenada por obra, o que permite cruzar com a mancha do INCT nível a nível, a mesma ligação que a Tela 1 usa. E a pergunta tem uso prático para o parceiro e para a população: mostra quanto da verba da reconstrução está em lugares que alagam, e quantos equipamentos públicos, como unidades de saúde, escolas e centros esportivos, ficam expostos já em cheias menores que a de 2024. A soma dos valores dessas obras é chamada de \textbf{investimento exposto}; não é custo de dano.

O painel financeiro da mesma fonte dá o contexto da cidade inteira: quanto a Prefeitura estima precisar, quanto pediu à União e quanto recebeu. O Copernicus entra como referência da enchente real. O TCE-RS serviu só para conferir os valores das obras.

\section{Uso de ferramenta de IA}

Utilizei o assistente Claude Code (Anthropic) no levantamento e na coleta das bases, na verificação do \campo{robots.txt} e das licenças e na redação deste relatório. As escolhas e as conclusões são minhas, e conferi as fontes.

\end{document}
